\cotefile{29}
\scanpage{21}
% archive sheet number: 20
% Formula spaces below reproduce unfilled spaces in the typescript; no missing signs have been supplied.
\deleted{$(a_i)_{i\quad I}$ de $\underline{O}_S$, \unsure{un|une}} une famille $(D_i)_{i\quad I}$ de diviseurs (de Cartier) sur $S$. \unsure{Bour|Pour} tout $S'$ étale sur $S$, on désigne par $\underline{C}_{S'}$ la catégorie suivante: les objets sont les \deleted{\unsure{couples triples}} \inserted{quadruples $M=$} $(\underline{a},\underline{n},\inserted{G,}\underline{u})$, où $\underline{a}=(a_i)_{i\quad I}$ est une famille d'équations pour les $D_{iS'}$, $\underline{n}=(n_i)_{i\quad I}$ une famille d'entiers $1$ premiers aux caractéristiques résiduelles de $S'$, \deleted{\unsure{et} $u:\unsure{G}\quad{}_{\underline{n}}$} $G$ un groupe fini, et $u:G_{S'}\quad{}_{\underline{n}}S'$ un isomorphisme de schémas en groupes sur $S'$. 2tant donné un deuxième tel quadruplet \inserted{$M'=$} $(\underline{a}',\underline{n}',G',u')$, \deleted{un morphisme du second dans le premier est $M'$ dans $M$} $\operatorname{Hom}_{S'}(M',M)$ est vide sauf si $\underline{n}'\quad\underline{n}$, auquel cas c'est l'ensemble des $\underline{b}=(b_i)_{i\quad I}$ tels que $\underline{b}^{n}=\underline{a}'\underline{a}^{-1}$. La composition des morphismes,et la définition des foncteurs changements de base pour une flèche $S''\quad S'$ dans $\underline{\mathrm{Et}}_S$, sont clairs. Le foncteur cartésien $\underline{C}\quad(\mathrm{Sch})/S$ est défini en associant à $M=(\underline{a},\underline{n},G,\underline{u})$ le revêtement Kummérien élémentaire $Z_{\underline{a}}^{\underline{n}}$, \deleted{et à une} \deleted{\unsure{flê|flè}} considéré comme un revêtement à groupe d'opérateurs $G$, grâce à l'isomorphisme $G\quad{}_{\underline{n}}$. On définit de faç\deleted{\unsure{x}}on évidente la valeur de $\underline{r}$ sur les flèches de $\underline{C}_{S'}$, et sa compatibilité avec les changements de base. On voit alors qu'on est sous les conditions de l'exemple 3. De plus, le foncteur $\underline{r}$ est fidèle, et il est pleinement fidèle lorsque $S$ est \deleted{noethérien aux points de} localement noethérien, et de plus normal en les points des $\operatorname{supp}D_i$, et si de plus les $D_i$ sont réduits, et toute \inserted{composante irréductible \unsure{de'un|d'un} $D_i$ est distincte de toute comp. d'un $D_{\unsure{i|j}}$ pour $j\ne i$}

b) Supposons donné un diviseur $D$ sur $S$, \deleted{une famille couvrante $\unsure{S'|S}\quad S$} et considérons le faisceau \unsure{éur|śur|sur} $\underline{\mathrm{Et}}_S$ associé au préfaisceau qui, à tout $S'$, associe l'ensemble des \deleted{diviseurs} ensembles \deleted{$\unsure{\underline{Q}|\underline{\Omega}}$} \inserted{$(D_i)$} $\unsure{\theta|0}$ de diviseurs sur $S'$ ayant les propriétés suivantes: $D_{\unsure{S'|S''}}=\quad D_i$, et pour tout $s'\quad S'$, les restrictions \deleted{non nulles} à $\text{Sêc}(\underline{O}_{S',s'})$ des $D_i$ qui sont non nulles, sont deux à deux disjointes. Soit donnée une section du
\scanpage{22}
% archive sheet number: 21
faisceau associé au préfaisceau précédent (NB la condition canularesque mise sur les $(D_i)$ assure que l'on a bien un préfaisceau). Alors localement pour la topologie étale, on a donc une situation $(D_i)$ comme dans a). Cela suffit pour définir une donnée de ramification $\underline{R}$ sur $\unsure{B|S}$, généralisant d'ailleurs celle envisagée dans a) (du moins lorsque la famille $(D_i)_{i\quad I}$ de a) satisfait à la condition supplémentaire locale qu'on\handwritten{\unsure{g|?}} vient de formuler). Je vois deux cas particuliers ayant de l'intérêt\unsure{;.|;|:} \deleted{\unsure{fi|si}} et \deleted{1\textsuperscript{\uline{o}})} où $\quad$ peut être défini en termes d'un diviseur $D$ donné sur $S$:

1\textsuperscript{\uline{o}}) $S$ est localement noethérien, $D$ est à croisements normaux, on prend comme $D_i$ les composantes irréductibles locales régulières de $D$ (qui ne sont définies que localement pour la topologie étale).

2\textsuperscript{\uline{o}}) $S$ est au dessus d'un préschéma $T$, $D$ est à croisements normaux relativement à $T$, on prend comme $D_i$ les composantes locales lisses sur $\unsure{\mathbb{S}|S|8}$ (qui ne sont définies que localement pour la topologie étale).

Dans le cas 1\textsuperscript{\uline{o}}), la catégorie $\operatorname{Rev}(_{\underline{R}}(S)$ est équivalente à une sous-catégorie pleine de la catégorie des revêtements de $S$, et sauf erreur (il faudrait le vérifier) il en est de même dans le cas 2\textsuperscript{\uline{o}}). Les revêtements en question s'appellent, dans le\deleted{\unsure{ur}} cas 1\textsuperscript{\uline{o}}), revêtements modérément ramifiés au sens large relativement au diviseur $D$ (NB on vérifie en effet que c'est une notion plus large que celle de ``revêtement modérément ramifié par rapport à $D$'', (signifiant : par rapport à la famille réduite à $D$). Le théorème d'Abhyankar permet de caractériser ces revêtements par la condition d'être étales sur \deleted{$S\quad\unsure{\underline{a}}$} $S-\operatorname{supp}D$, normaux en les points au dessus de $\operatorname{supp}D$, toute composante irréductible en \unsure{dómine|dômine|domine} une de $S$, enfin ils sont modérément ramifiés au sens classique en le\inserted{s} points maximaux de $D$. Dans le cas 2\textsuperscript{\uline{o}}, il y aurait lieu d'appeler les revêtements obtenus ``modérément ramifiés pour $D$, relativement à la
\scanpage{23}
% archive sheet number: 22
base $T$ ''. Lorsque $T$ est régulier, les notions 1\textsuperscript{\uline{o}} et 2\textsuperscript{\uline{o}} coïncident. Le théorème de spécialisation pour le groupe fondamental tame peut se démontrer, par la technique de SGA 1, quand on utilise la notion de revêtement modérément ramifié au sens large, relativement aux\unsure{diviseurs|diviseur} $D_{\bar{t}}$ \inserted{dans les $S_{\bar{t}}$} ($D$ étant à croisement normaux relativement à $T$, $S$ propre et lisse sur $T$).

\uline{Remarques} Evidemment, pour tout $S'$ objet de $\underline{\mathrm{Et}}_S$, on définit de façon évidente $\operatorname{Rev}_{\deleted{\underline{R}}}^{\inserted{\underline{R}}}(S')$, en considérant la restriction de $\underline{R}$ à $\underline{\mathrm{Et}}_{S'}\quad(\underline{\mathrm{Et}}_S)/S'$. Pour $S'$ variable, on obtient ainsi une catégorie fibrée $\underline{\mathrm{Rev}}^{\underline{R}}$ sur $\underline{\mathrm{Et}}_S$, dont les fibres $\underline{\mathrm{Rev}}_{S'}^{\underline{R}}$ sont les $\operatorname{Rev}^{\underline{R}}(S')$, et un foncteur cartésien
\[
\underline{r}:\quad\underline{\mathrm{Rev}}^{\underline{R}}\qquad\underline{\mathrm{Rev}},
\tag*{(*)}
\]
où on désigne pour abréger par $\underline{\mathrm{Rev}}$ la catégorie \inserted{fibrée sur $\underline{\mathrm{Et}}_S$} des revêtements des différents objets $S'$ de $\underline{\mathrm{Et}}_S$. (NB on suppose que pour tout $S'$, le foncteur réalisation géométrique sur $\underline{C}_{S'}$ est à valeurs dans les revêtements de $S'$, et que le foncteur réalisation géométrique sur $\operatorname{Rev}^{\underline{R}}(S')$ est bien défini). Il faut insister sur le fait que pour nous, la ``donnée de ramification'' $\underline{R}$ est un objet de nature essentiellement auxiliaire, destiné à définir la catégorie fibrée $\underline{\mathrm{Rev}}^{\underline{R}}$ (qui est même un \uline{champ}) et le foncteur cartésien $\underline{r}$ de (*); c'est (*) qui est l'objet véritable de l'étude. On regardera donc deux données de ramification $\underline{R}$, $\underline{R}'$ comme essentiellement équivalentes, si elles définissent des foncteurs (*) équivalents (par une équivalence, bien sûr, induisant l'identité sur $\underline{\mathrm{Rev}}$). On peut se convaincre sur des exemples que $\underline{R}$ et $\underline{R}'$ peuvent avoir des allures très différentes. Ainsi, sous les conditions de l'exemple 3), si $\underline{R}'$ est défini par une sous-catégorie fibrée \deleted{de} $\underline{C'}$ de $\underline{C}$
\scanpage{24}
% archive sheet number: 23
telle que pour tout objet $M$ d'un $\underline{C}_{S'}$, il existe un recouvrement de $S'$ par des $S'_i$ tel que pour tout $i$, $M_{S'_i}$ soit \deleted{dominé} majoré par un $M_i$ objet de $\underline{C}'_{S'_i}$, \deleted{alors} et supposant que le foncteur réalisation géométrique pour $\underline{R}'$ est celui induit pour le foncteur réal. géom. pour $\underline{R}$, - alors on constate que $\underline{R}$ et $\underline{R}'$ sont ``équivalents''. Comme cas particulier, signalons l'exemple de la donnée kummérienne a), dans le cas où les $D_i$ admettent des équations globales par des sections $c_i$ de $\underline{O}_S$. Alors on peut considérer le système projectif des $Z_{\underline{c}}^{\underline{n}}$, pour $\underline{n}=(n_i)_{i\quad I}$ variables, premiers \deleted{aux \unsure{caractéristiques|caractéritiques} résiduelles} à la caractéristique de $S$ (supposée exister). Considérons de plus le système projectif des groupes $\underline{Z}_{\underline{n}}=\quad\underline{Z}/n_i\underline{Z}$, et supposons donné un isomorphisme entre celui-ci et le système projectif des ${}_{\underline{n}}$ (il en existe lorsque $S$ est un schéma sur un corps algébriquement clos, par exemple). On peut alors considérer les $Z_{\underline{c}}^{\underline{n}}$ comme des schémas à groupes $\underline{Z}_{\underline{n}}$, formant un système projectif, \deleted{d'où} donc on est sous les conditions de l'exemple 2, définissant une donnée de ramification $\underline{R}'$ sur $S$. On \unsure{constat|constate} alors que cette donnée est équivalente à la donnée $\underline{R}$ décrite dans le cas particulier kummérien a). Comme autre exemple de données de ramification équivalentes\deleted{\unsure{:|;}}, signalons le suivant, où $\underline{R}$ est la donnée de ramification considérée dans l'exemple 1. On se donne un revêtement principal $P$ de $S$, de groupe ${\unsure{H|R|G}}'$ (opérant à droite), et on suppose que $G$ opère dessus librement à gauche, en commutant aux opérations de $G'$. Soit alors $Z'=Z\times^{G}P$, considéré comme revêtement à groupe d'opérateurs $G'$, d'où une donnée de ramification $\underline{R}'$. Cette dernière est équivalente à $\underline{R}$, comme on voit en considérant l\deleted{es}'équivalences \deleted{quasi-inverses l'une} de \deleted{l'autre} \deleted{entre} $\underline{\mathrm{Rev}}^{\underline{R}}$ \deleted{et}\inserted{dans} ${\unsure{\underline{\mathrm{Rev}}|\mathrm{Rev}}}^{R'}$ opérateurs \illegible{3}
\scanpage{25}
% archive sheet number: 24
évidente lorsque $P=G'_S$ avec l'opération régulière droite de $G'$, et l'opération gauche de $G$ est donnée par un monomorphisme $G\quad G'$; on se ramène à ce cas particulier par localisation étale.

On peut se proposer de normaliser la notion de donnée de ramification, pour éliminer l'arbitraire dans le choix d'un $\underline{R}$. Mais alors \unsure{bl|il|dl} semble qu'on soit infailliblement conduit à prendre comme catégorie $\underline{C}$ la catégorie $\underline{\mathrm{Rev}}^{\underline{R}}$ elle-même qu'on se propose \unsure{dé|de} construire et d'étudier ! (Comparer avec la situation analogue pour les relations entre sites et topos, l'objet auxiliaire étant le site, l'objet intrinsèque de nature géométriquenétant le topos associé à ce site).

2. \uline{Propriétés de la catégorie} $\operatorname{Rev}^{\underline{R}}(S)$. \uline{Le groupe fondamental} ${}^{\underline{R}}_1(S,\quad)$.

On constate immédiatement que la catégorie $\operatorname{Rev}^{\underline{R}}(S)$ hérite de toutes les propriétés d'exactitude \deleted{déjà} établies pour la catégorie des revêtementsétales d'un schémas donné, et dont une partie sont énoncés dans Exp V n\textsuperscript{\uline{o}}4 (conditions G1 à G3). De façon précise, on peut dire que dans la catégorie $\operatorname{Rev}^{\underline{R}}(S)$, les limites projectives finies et les limites inductives finies existent, et satisfont aux mêmes propriétés d'exactitude que dans la catégorie des ensembles (telle l'exactitude du foncteur changement de base dans cette catégorie); \deleted{de plus;} \inserted{en particulier} tout morphisme de la catégorie se factorise en le produit d'un épimorphisme et d'un monomorphisme, \inserted{les épimorphismes et monomorphismes sont \deleted{\unsure{stri}} effectifs; enfin} tout monomorphisme est un isomorphisme avec un sommande direct. On fera attention que le foncteur ``réalisation géométrique'' sur $\operatorname{Rev}^{\underline{R}}(S)$ commute bien aux limites inductives finies (et \deleted{par suite aux} en particulier aux sommes finies), mais non pas en général aux limites projectives finies, et en particulier, pas aux produits de deux objets.

Pour appliquer la théorie de Galois générale de V 4, il faut dis%
\scanpage{26}%
% archive sheet number: 25
poser de plus d'un \uline{foncteur fibre}
\[
\unsure{F|\mathbb{F}}:\operatorname{Rev}^{\underline{R}}(S)\qquad(\mathrm{Ens\ finis}),
\]
i.e. d'un foncteur exact\inserted{(conditions G4,G5 de. V 4)}. Dans le cas de l'exemple 1, tout point géométrique de $Z$ fournit un tel foncteur fibre, et ce foncteur est conservatif (condition G6) si on suppose que toute partie ouverte et fermée de $Z$ stable par $G$ est vide ou égale à $Z$. De même, dans le cas de l'exemple 2, on obtient un \deleted{\unsure{par}} foncteur fibre en prenant un point géométrique de $Z=\lim Z_i$. Bour obtenir un foncteur fibre dans le cas plus général de l'exemple 3, on choisit un point géométrique $\qquad S$ de $S$ \deleted{tel que} qui n'appartienne pas à\deleted{\unsure{u}} ``l'ensemble de ramification''de $\underline{R}$, i.e. tel que pour tout $S'$ dans $\mathrm{Et}_S$ et tout objet $M$ de $\underline{C}_{S'}$, \unsure{le|l'} \deleted{\unsure{objet}} \inserted{revêtement à} \deleted{à} opérateurs $\underline{r}(M)$ \deleted{sur $S'$}\inserted{(groupe $G$)} de $S'$ soit principal en les points de $S'$au dessus du l'image \inserted{$\unsure{s|S}$} de $\quad$ dans $S$. Alors \deleted{pour tout revêtement étale} \deleted{$T$ \unsure{à groupe $G$} \unsure{de} \unsure{à l'ob|l'ob} d'un schéma à opérateurs $G$} on constate que pour tout objet $T$ de $\operatorname{Rev}^{\underline{R}}(S)$, $\underline{r}(T)$ est étale enr $S$ en $s$, et le foncteur $T\quad\underline{r}(T)(\quad)$ (fibre géométrique en $\quad$) est bien exact. Pour qu'il soit conservatif, il faut et il suffit que le site $\underline{C}$ (pour la topologie image inverse de celle de $\underline{\mathrm{Et}}_S$) soit connexe; \deleted{i.e. que tout crible couvrant de ce site} si on suppose, avec les notations de l'exemple 3, que pour tout $M$, définissant $(\underline{r}(M),G)$, $\underline{r}(M)/G\qquad S'$ soit un isomorphisme (ce qui sera sans doute le cas dans tous les cas intéressants), alors il revient au même de dire que $S$ est connexe. Appliquant la théorie de V 4, on trouve donc un groupe pro-fini, noté ${}^{\underline{R}}_1(S,\quad)$, \deleted{qui classifi} et une équivalence de la catégorie des ensembles finis où ce groupe opère continûment, avec la catégorie $\operatorname{Rev}^{\underline{R}}(S)$.

3. \uline{Fonctorialités pour} $S$ \inserted{\uline{et} $\underline{R}$} \uline{variables}

\unsure{Crnsidérons|Considérons} un morphisme $f:T\qquad S$ de schémas, d'où un foncteur
